\section{Smoothed weights and area calculation}\label{app:area}
We interpolate the weights $\exp(-\chi^2_{\rm int}/2)$ on a fixed triangulation of the 1285 sampled points, using the same geometry for every foreground treatment and redshift combination. The surface is evaluated on an $800\times800$ grid spanning $0.195\le\Omega_m\le0.425$ and $-1.325\le w\le-0.670$. Before Gaussian smoothing, missing weights outside the sampled convex hull are continued from the nearest sampled point, and the rectangular edges use a reflected continuation. These choices can affect the smoothed weights inside the hull and are part of the prescription.

We use the same Gaussian widths in parameter units for every surface, $\sigma_{\Omega_m}=0.0068836$ and $\sigma_w=0.0196033$. They originate from an empirical bandwidth based on 2.5 times the median nearest-neighbour distance in the unit-scaled scan rectangle, rather than from the covariance or an independent error calibration. This smoothing acts on the parameter weights, not on the measured genus curves or on the density field smoothed at $R=20\,h^{-1}\mathrm{Mpc}$.

Within the convex hull, weights are ordered from highest to lowest and the threshold is set where their cumulative sum reaches the chosen fraction of the total. The enclosed area is the number of selected grid elements times their area in $(\Omega_m,w)$. Thresholds are determined separately for each smoothed surface. For the joint surfaces, doubling the grid resolution at fixed parameter-unit kernel widths changes the areas by less than 0.6 per cent; this tests integration sampling, not the kernel choice or spatial convergence of the density field.
